\section{Discussion}

\subsection{Key Findings}

\textbf{mypy augmentation provides substantial but mechanism-agnostic benefit.}
Our primary finding is a +8.94pp pass@1 improvement from adding mypy to
execution-based repair loops on HumanEval+ (GPT-4o-mini, $k$=5, three seeds
consistent). The type-specificity hypothesis --- that mypy helps specifically
for type-error problems --- is not supported. Both conditions achieve identical
90.9\% repair rates on type-error problems (h-m2, $n$=22, differential=0.000).
The most parsimonious interpretation --- pending ablation confirmation (see L2)
--- is \emph{general diagnostic context enrichment}: mypy output adds structured
text to the repair prompt that may activate GPT-4o-mini's latent knowledge about
Python error patterns, improving repair quality across problem types. This
suggests that other structured diagnostic channels (linters, semantic analyzers)
might provide similar benefits, independent of whether they target the specific
failure mode.

\textbf{Benchmark-level signal divergence changes the applicability picture.}
The 70\% vs 0\% mypy error rate between HumanEval+ and MBPP+ is a new empirical
finding. HumanEval+ failures are dominated by undefined-name errors; MBPP+
failures appear to be pure algorithmic correctness errors invisible to mypy.
This divergence provides a practical guideline: a practitioner should first
profile the mypy error rate in failing solutions before adding mypy feedback ---
if it is low ($<$10\%), mypy will not activate frequently enough to drive
improvement.

\textbf{Z3 counterexample feedback: sound architecture, insufficient activation.}
The Z3 null result ($-$0.9pp) is informative: the spec generation pipeline
achieves 91.1\% validity, confirming LLM-generated Z3 specs are mostly correct
on arithmetic HumanEval+ problems. The bottleneck is CE activation (16.1\%).
Z3 counterexamples would be most valuable on problems where base pass@1 is
well below the ceiling ($<$60\%), leaving genuine headroom for guided repair.

\subsection{Limitations}

\textbf{L1: MBPP+ pass@1 comparison not completed.} All quantitative pass@1
claims (+8.94pp) are HumanEval+-specific. The h-m3 MBPP+ experiment was started
(122/378 tasks, seed=42, Condition A only) but not completed due to compute
constraints. Given the 0\% mypy error rate on MBPP+ (h-e1, h-m1), Condition B
is unlikely to show improvement on MBPP+ --- but this remains unconfirmed. We
restrict all pass@1 claims to HumanEval+.

\textbf{L2: Type-specificity mechanism alternative unverified.} We propose
general context enrichment as the active mechanism, but the ablation study
needed to confirm it (Condition B-Null: same prompt length as B, but with
generic filler text instead of mypy output) was not executed. Without this
ablation, we cannot distinguish mypy \emph{content} from mypy \emph{length} as
the active ingredient. This is the highest-priority future experiment.

\textbf{L3: Single model (GPT-4o-mini).} The repair benefit may be
model-capability-dependent. Replication with at least one open-source code
model is straightforward future work.

\textbf{L4: Z3 feedback at high base pass@1.} The arithmetic subset experiment
was conducted at base pass@1\_B = 86.6\%, leaving 13.4\% headroom. The null
result may not generalize to settings with lower base pass@1.

\subsection{Broader Impact}

This work establishes that adding mypy to LLM repair loops provides substantial,
reproducible improvement in Python code generation quality. Practical deployment
is low-cost (one subprocess call per repair round, ${\sim}$30ms latency). The
finding that the improvement is not type-specific suggests that other structured
feedback channels deserve controlled experimental evaluation rather than being
bundled with execution feedback and assumed beneficial. The Z3 pipeline
(91.1\% valid specs) establishes feasibility for automated formal specification
generation, a building block for future formal verification of LLM-generated
code. We do not foresee significant negative impacts: improving automated repair
quality does not introduce new failure modes beyond those of LLM-generated code
generally.
